\documentclass[11pt]{article}
\usepackage[margin=0.94in]{geometry}
\usepackage{amsmath,amssymb,amsthm,booktabs}
\usepackage[hidelinks]{hyperref}
\hypersetup{pdfauthor={Ian Walker},pdftitle={The Exact Positive-Definiteness Dimension of the Riemann Xi Function}}
\newtheorem{theorem}{Theorem}
\newtheorem{lemma}{Lemma}
\newcommand{\R}{\mathbb R}
\newcommand{\ee}{\mathrm e}
\title{The Exact Positive-Definiteness Dimension\\of the Riemann Xi Function}
\author{Ian Walker}
\date{Referee draft, November 25, 2026}

\begin{document}
\maketitle
\begin{abstract}
The classical radial Fourier dimension walk reduces the Euclidean
positive-definiteness problem for the normalized Riemann Xi function to signs
of derivatives of its theta density. Monotonicity due to Wintner and
strict convexity due to Csordas give positivity through dimension five.
We prove that the nine-dimensional density is positive at every nonzero
radius, giving strict positive definiteness through dimension nine.
In dimension ten an explicit finite Gram form has a negative value,
certified by grouping its terms into 739 distance shells. Thus nine is
the exact maximal dimension. Lean verifies the analytic lower bound and
the finite-witness algebra. Arb encloses the transcendental evaluations
for the upper bound.
\end{abstract}

\section{Statement and context}
Write
\[
\xi(s)=\tfrac12s(s-1)\pi^{-s/2}\Gamma(s/2)\zeta(s),
\qquad \Xi(t)=\xi(\tfrac12+it),
\qquad F(t)=\Xi(t)/\Xi(0),
\]
with the removable singularities filled in. The functional equation and
conjugation make $F$ real and even on $\R$. A kernel is positive
semidefinite if every finite Gram form is nonnegative. Strict positive
definiteness requires the form to be positive whenever the points are
distinct and the coefficient vector is nonzero.

Schoenberg's radial classes are nested under restriction to Euclidean
subspaces \cite{Schoenberg}. To determine where $F$ leaves this hierarchy,
let $f_d(r)$ denote its radial Fourier density in dimension $d$, and set
$q(s)=f_1(\sqrt{s})$. With the angular-frequency convention of Section~2,
the classical dimension walk \cite{GT} gives, for $s>0$,
\[
 f_{2j+1}(\sqrt{s})=(-1/\pi)^j q^{(j)}(s).
\]
Thus each odd-dimensional test asks for one derivative sign. The
one-dimensional density is the positive classical theta kernel up to
normalization \cite{BPY,Nakamura}. In the notation of Csordas
\cite[(2.1)--(2.2)]{Csordas98}, it gives
$q(s)=\Phi_{\rm C}(\sqrt{s}/2)/Z$, with $Z>0$ as in Section~2.
Wintner's monotonicity \cite{Wintner} gives $q'<0$, and strict convexity
\cite[Theorem 2.12]{Csordas98} gives $q''>0$. Hence $f_3$ and $f_5$
are positive away from zero. The strict Gram argument in Section~3
and restriction already yield strict positive definiteness through
dimension five.

The step proved here is $q^{(4)}(s)>0$ for every $s>0$. It gives
$f_9>0$ away from zero. Its two-coordinate marginal gives $f_7>0$,
and restriction supplies dimensions six and eight. In dimension ten
a single explicit negative Gram form gives the matching upper bound.

\begin{theorem}\label{thm:main}
For an integer $d\ge1$, the kernel $K_d(x,y)=F(\lVert x-y\rVert)$
on $\R^d$ is positive semidefinite if and only if $d\le9$.
For $d\le9$ it is strictly positive definite. For $d=10$ the explicit
witness in Section~\ref{sec:witness} has normalized Gram form
strictly less than $-0.01259636712$.
\end{theorem}

Sra proves that for each $\alpha>0$ the real-axis reciprocal profile
$[\xi(1/2)/\xi(1/2+\lVert u\rVert)]^\alpha$ is strictly positive definite
in every Euclidean dimension \cite[Theorem 1.1, (1.4)]{Sra}. Our profile
$F(\lVert u\rVert)=\xi(1/2+i\lVert u\rVert)/\xi(1/2)$ evaluates $\xi$
on the critical line. Its positive-radius zeros already rule out
positive definiteness in all dimensions simultaneously: Schoenberg's
all-dimensional representation would express $F(r)$ as a probability
mixture of $\exp(-tr^2)$, strictly positive at every finite $r$.
Determining the last admissible dimension requires the two sign results
above.

Other related work includes Romik's three-dimensional self-transform of
an auxiliary function \cite{Romik}, the van Dantzig pair studied in
\cite{KPS}, the reciprocal transforms in \cite{Grochenig}, and the
smoothed zero sums in \cite{SraSmooth}.

\section{Theta normalization and radial lifting}
For $r\in\R$ set
\[
 X_n(r)=\pi n^2\ee^{2r},\quad
 g(r)=\sum_{n\ge1}\ee^{r/2-X_n(r)},\quad
 \Phi(r)=\sum_{n\ge1}\ee^{r/2-X_n(r)}X_n(r)(2X_n(r)-3).
\]
The series and their derivatives converge locally uniformly by Gaussian
majorants. Each derivative is a sum with the same exponential factor
times a polynomial in $X_n$. These sums decay faster than any exponential
as $r\to+\infty$. For example, after absorbing the polynomial into
$\ee^{X_n/2}$, split $\ee^{-X_n/2}$ into two factors, one bounded by
$\ee^{-\pi n^2/4}$ on $r\ge0$ and the other by
$\ee^{-\pi\ee^{2r}/4}$.
Jacobi inversion makes $H(r)=\ee^{r/2}+2g(r)$ even. Thus
$g'(0)=-1/4$. Termwise differentiation gives
$g''-g/4=2\Phi$ and $H''-H/4=4\Phi$, so $\Phi$ is smooth, even and
Schwartz. It is strictly positive on $[0,\infty)$ since $X_n(r)\ge\pi$.

Splitting the theta Mellin integral at one and applying inversion to its
lower half gives, for $s=1/2+it$,
\[
 \Xi(t)=\tfrac12-2(t^2+\tfrac14)\int_0^\infty g(r)\cos(tr)\,dr.
\]
Indeed, for $\psi(x)=\sum_{n\ge1}\ee^{-\pi n^2x}$ the pole-subtracted
completed zeta is the integral on $[1,\infty)$ of
$\psi(x)(x^{s/2}+x^{(1-s)/2})\,dx/x$. The substitution $x=\ee^{2r}$
gives $4g(r)\cos(tr)\,dr$, and the pole terms contribute $1/2$ to $\xi$.
Two integrations by parts using $g'(0)=-1/4$ therefore yield
\[
 \Xi(t)=4\int_0^\infty \Phi(r)\cos(tr)\,dr,
 \qquad
 F(t)=\int_{\R}\ee^{itr}\frac{\Phi(r)}{Z}\,dr,
 \qquad Z=\int_{\R}\Phi(r)\,dr>0.
\]

We use angular frequencies, so a spectral density $f_d$ is normalized
by $F(|x|)=\int_{\R^d}\ee^{ik\cdot x}f_d(k)\,dk$.
The even Schwartz profile $F$ gives a Schwartz radial function in each
dimension. Let $f_d(r)$ denote its real radial Fourier density. Fourier
uniqueness, restriction to $\R^d$, and polar integration in the last two
coordinates give
\[
 f_d(r)=2\pi\int_r^\infty s f_{d+2}(s)\,ds,
 \qquad f_{d+2}(r)=-\frac{f_d'(r)}{2\pi r}\quad(r>0).
\]
Starting with $f_1(r)=\Phi(r)/Z$ and iterating four times yields
\begin{equation}\label{eq:d9}
 f_9(r)=\frac{N(r)}{(2\pi)^4Zr^7},\qquad
 N(r)=r^3\Phi^{(4)}(r)-6r^2\Phi^{(3)}(r)+15r\Phi''(r)-15\Phi'(r).
\end{equation}
Only $r>0$ is needed for positivity almost everywhere. The Lean proof
implements the marginal identity directly and propagates integrability
by positivity and Fubini.

\section{An analytic sign certificate in dimension nine}
Let $P_0(x)=2x^2-3x$ and define
\begin{equation}\label{eq:recurrence}
 P_{j+1}(x)=(\tfrac12-2x)P_j(x)+2xP_j'(x).
\end{equation}
Then $\Phi^{(j)}(r)=\sum_{n\ge1}\ee^{r/2-X_n(r)}P_j(X_n(r))$.

\paragraph{Near the origin.}
For $0\le r\le1/5$, put $x=X_1(r)$. Elementary bounds give
$3\le x\le47/10$. In fact $3.1415<\pi<3.1416$, and the exponential
series bounds $\ee^{2/5}$ above by $1.492$ and below by its degree-five
Taylor sum. The upper bound follows by bounding the terms beyond
degree five by a geometric series of ratio at most $(2/5)/7$.
Define
\[
 L(x)=8000\bigl(1+3(x-3)+\tfrac12(3(x-3))^2\bigr).
\]
On $x\ge3$, $L(x)\le\ee^{3x}$ by the Taylor lower bounds
$\sum_{k=0}^{19}9^k/k!>8000$ and $\ee^y\ge1+y+y^2/2$ for $y\ge0$.
The first two terms of $\Phi^{(8)}$ have positive sum since
$P_8(x)>0$ and $L(x)P_8(x)+P_8(4x)>0$ on this interval.
For $n\ge3$, $X_n(r)\ge27$, so $u=X_n(r)-19\ge8$.
Every coefficient of $P_8(19+u)$ is positive, hence so are the remaining terms.
The exact polynomial checks are given below.

Put $p_j=\Phi^{(j)}$ and define
\[
 A(r)=rp_6(r)-p_5(r),\qquad
 B(r)=r^2p_5(r)-3rp_4(r)+3p_3(r).
\]
Evenness gives $p_7(0)=A(0)=B(0)=N(0)=0$. The identities
$p_7'=p_8$, $A'=rp_7$, $B'=rA$, and $N'=rB$ show successively
that each of $p_7,A,B,N$ is positive on $(0,1/5]$.

\paragraph{All larger radii.}
Write $r=1/5+z$, with $z\ge0$, and
$b(z)=(937/200)(1+2z+2z^2)$.
The degree-five Taylor lower bound for $\ee^{2/5}$ and the lower
bound $\pi>3.1415$ give $\pi\ee^{2/5}\ge937/200$.
Thus $X_n(r)\ge b(z)$ for every $n\ge1$. The polynomial
\[
 T(r,x)=r^3P_4(x)-6r^2P_3(x)+15rP_2(x)-15P_1(x)
\]
has positive constant coefficient and nonnegative coefficients after
the substitutions $r=1/5+z$ and $x=b(z)+u$. Hence every summand
of $N(r)$ is strictly positive for $r\ge1/5$.

\paragraph{Exact arithmetic checks.}
For a polynomial $q$ on $[3,47/10]$, let
$q(3+17t/10)=\sum_{j=0}^d a_j t^j$. Its degree-$d$ Bernstein
coefficients are $\beta_i=\sum_{j=0}^i a_j\binom{i}{j}/\binom{d}{j}$.
The basis functions are nonnegative and sum to one on $[0,1]$.
The following minima follow by exact rational arithmetic from
\eqref{eq:recurrence}. The supplement includes a short independent
checker for these values.

\smallskip
\begin{center}
\small
\begin{tabular}{lll}
\toprule
Polynomial & Basis & Minimum coefficient \\
\midrule
$P_8(x)$ & Bernstein, degree 10 & $5029330761/256$ \\
$L(x)P_8(x)+P_8(4x)$ & Bernstein, degree 12 & $10192605107247/64$ \\
$P_8(19+u)$ & Monomials in $u$ & $512$ \\
$T(1/5+z,b(z)+u)$ & Nonzero monomials in $z,u$ & $32/125$ \\
\bottomrule
\end{tabular}
\end{center}
The constant coefficient of the last polynomial is
$1328632816292109/250000000000000>0$.

\paragraph{Strictness.}
Equation \eqref{eq:d9} gives $f_9>0$ almost everywhere. For
distinct $x_i$ and a nonzero coefficient vector $c$, the Gram form is
\[
 \int_{\R^9} f_9(k)
 \left|\sum_i c_i\ee^{ik\cdot x_i}\right|^2\,dk.
\]
The exponential sum is continuous and cannot vanish identically.
For example, choose a line direction outside the finitely many
hyperplanes perpendicular to pairwise differences of the $x_i$.
The resulting frequencies are distinct, and successive derivatives
at zero give an invertible Vandermonde system. The integral is
strictly positive. Restriction along isometric embeddings proves the
claim in all lower dimensions. Lean uses linear independence of
characters for the same strictness step.


\section{A finite witness in dimension ten}\label{sec:witness}
For $a\in\{0,\ldots,9\}^{10}$ use the point $x_a=7a$ and coefficient
$c_a=\prod_{\ell=1}^{10} w_{a_\ell}$, with
\[
 (w_0,\ldots,w_9)=(282,541,756,910,990,990,910,756,541,282).
\]
Define $P(z)=\sum_{i,j=0}^9 w_iw_j z^{(i-j)^2}$, $C_k=[z^k]P(z)^{10}$,
and $D=C_0=(\sum_iw_i^2)^{10}$. The normalized quadratic form is
\[
 \frac1D\sum_{a,b}c_ac_bK_{10}(x_a,x_b)
 =\sum_{k=0}^{810}\frac{C_k}{D}\,F(7\sqrt{k}).
\]
This follows by multiplying the ten ordered-pair histograms. Exactly
739 coefficients are nonzero, including the diagonal. The normalization
$D$ is a positive integer.

Arb/Acb \cite{Johansson} encloses the factors in the defining expression
for $\xi$ and verifies $\xi(1/2)>0$. Each normalized shell contribution
is rounded outwards to integer multiples of $10^{-18}$. Both the
192-bit and 256-bit runs give the following enclosure for the sum
of these rounded intervals:
\[
 \frac{-12596367123836528}{10^{18}}
 \ \le\ \frac1D\sum_{a,b}c_ac_bK_{10}(x_a,x_b)
 \ \le\ \frac{-12596367123835788}{10^{18}}.
\]
The upper endpoint is strictly below $-0.01259636712$.
All nonzero shells are evaluated. The cutoff $293$ used to store
the rows in two lists only partitions the finite computation.
Embedding the grid into any higher Euclidean space preserves the Gram
form. This completes the proof of Theorem~\ref{thm:main}.

\section{Verification and reproducibility}
The supplementary sources are available at
\url{https://github.com/iqwalker-debug/xi-positive-definiteness-}.
The public snapshot is pinned to
\href{https://github.com/iqwalker-debug/xi-positive-definiteness-/commit/b85224897e27d51141b85c98bad5173c78211dce}{commit \texttt{b852248}},
with \href{https://github.com/iqwalker-debug/xi-positive-definiteness-/actions/runs/45646086586}{verification run \texttt{45646086586}}.
The Lean declarations \texttt{strict\_d9} and
\texttt{strict\_cutoff\_of\_arb} are in \texttt{XiStrict.lean}.
The D9 proof has no numerical assumptions. For D10, Lean proves the
obstruction from the row enclosures. Arb supplies those enclosures
externally. The supplement contains the full shell table, a second
aggregate finite-sum replay, the D9 interval cross-check, and exact
commands. Neither the analytic argument nor the witness uses
information about the zeros of zeta.


\section*{Acknowledgments}
I used Claude (Anthropic) and ChatGPT (OpenAI) in preparing this work.

\begin{thebibliography}{10}
\bibitem{Schoenberg} I.~J. Schoenberg, \emph{Metric spaces and positive
definite functions}, Trans. Amer. Math. Soc. 44 (1938), 522--536.
\href{https://doi.org/10.1090/S0002-9947-1938-1501980-0}{doi:10.1090/S0002-9947-1938-1501980-0}.
\bibitem{Wintner} A. Wintner, \emph{A note on the Riemann $\xi$-function},
J. London Math. Soc. 10 (1935), 82--83. See also \cite[Theorem 2.1(v)]{Csordas98}.
\bibitem{Csordas98} G. Csordas, \emph{Convexity and the Riemann
$\zeta$-function}, Glas. Mat. Ser. III 33 (1998), 37--50.
\href{https://web.math.pmf.unizg.hr/glasnik/33.1/33105.pdf}{Full text}.
\bibitem{BPY} P. Biane, J. Pitman and M. Yor, \emph{Probability laws
related to the Jacobi theta and Riemann zeta functions, and Brownian
excursions}, Bull. Amer. Math. Soc. 38 (2001), 435--465.
\href{https://doi.org/10.1090/S0273-0979-01-00912-0}{doi:10.1090/S0273-0979-01-00912-0}.
\bibitem{Nakamura} T. Nakamura, \emph{A complete Riemann zeta
distribution and the Riemann hypothesis}, Bernoulli 21 (2015), 604--617.
\href{https://doi.org/10.3150/13-BEJ581}{doi:10.3150/13-BEJ581}.
\bibitem{KPS} T. Konstantopoulos, P. Patie and R. Sarkar,
\emph{A new class of solutions to the van Dantzig problem, the Lee--Yang
property, and the Riemann hypothesis}, Ann. Inst. Fourier 74 (2024), 377--421.
\href{https://aif.centre-mersenne.org/item/?id=AIF_2024__74_1_377_0}{Article}.
\bibitem{GT} L. Grafakos and G. Teschl, \emph{On Fourier transforms of
radial functions and distributions}, J. Fourier Anal. Appl. 19 (2013),
167--179. \href{https://doi.org/10.1007/s00041-012-9242-5}{doi:10.1007/s00041-012-9242-5}.
\bibitem{Romik} D. Romik, \emph{Orthogonal polynomial expansions for the
Riemann xi function in the Hermite, Meixner--Pollaczek, and continuous Hahn
bases}, Acta Arith. 200 (2021), 259--329.
\href{https://arxiv.org/abs/1902.06330}{arXiv:1902.06330}.
\bibitem{Grochenig} K. Gr\"ochenig, \emph{Schoenberg's theory of
totally positive functions and the Riemann zeta function}.
\href{https://arxiv.org/abs/2007.12889}{arXiv:2007.12889}.
\bibitem{Sra} S. Sra, \emph{On some positivity properties of the Riemann
$\Xi$ function}, preprint (2026).
\href{https://arxiv.org/abs/2611.02944}{arXiv:2611.02944}.
\bibitem{SraSmooth} S. Sra, \emph{Some positivity properties of Riemann
$\Xi$ under smoothing}, unpublished preprint (2026).
\href{https://hexagonmath.org/2611.00345v1}{hexagon:2611.00345v1}.
\bibitem{Johansson} F. Johansson, \emph{Arb: Efficient arbitrary-precision
midpoint-radius interval arithmetic}, IEEE Trans. Comput. 66 (2017),
1281--1292. \href{https://doi.org/10.1109/TC.2017.2690633}{doi:10.1109/TC.2017.2690633}.
\end{thebibliography}
\end{document}
